\begin{abstract}
Blind image deblurring under extreme degradation fails for two
fundamental reasons: deterministic regression networks overfit to
training blur distributions and collapse beyond their training range,
while flow-based generative methods assume either a reliable posterior
mean or a known degradation operator---both of which break down when
blur is severe.
%
We identify the root cause geometrically: under extreme blur, the
straight-line trajectory of standard Rectified Flow (RF) traverses
low-probability regions of the clean image manifold, producing
artefacts or semantic incoherence.
%
We propose \textbf{SemFlow}, a framework that reshapes the RF velocity
field using spatially grounded semantic tokens extracted by a
Joint Embedding Predictive Architecture (JEPA) encoder.
%
Unlike CLIP-based conditioning, JEPA is pre-trained to predict
abstract representations from masked and corrupted inputs, making its
context encoder structurally robust to the extreme degradation it must
encode.
%
Multi-scale JEPA patch tokens are injected into the velocity field via
time-varying cross-attention, steering the flow trajectory toward
semantically consistent sharp-image regions without any blur kernel
information.
%
An auxiliary JEPA self-consistency loss further constrains restored
outputs to match ground-truth abstract structure, complementing
pixel-level and perceptual objectives.
%
We define an operational extreme-blur regime (effective kernel diameter
$\geq\!33$\,px), introduce the \textbf{Extreme-GoPro-5$\times$}
benchmark constructed with physics-based synthetic blur, and demonstrate
through ablation that trajectory curvature induced by JEPA conditioning
is the decisive factor distinguishing SemFlow from prior flow-based
and VLM-conditioned baselines.
\end{abstract}
